\section{A positive renewal equation retains the arithmetic forcing (NS-68)}
\label{sec:ns68-renewal}

\paragraph{Status.}
The first multiplicative interval of the canonical response gives a
particularly simple signed arithmetic quantity. Its exact renewal
operator is positive and contractive in every strictly positive power
weight. This controls inversion of the operator but does not improve
the growth exponent of its forcing. The required arithmetic estimate
is still RH-equivalent. The argument below is a direct calculation;
no claim of priority for these reformulations is made.

For real $x>0$ put
\[
 M(x)=\sum_{n\leq x}\mu(n),\quad s(x)=\sum_{n\leq x}\frac{\mu(n)}n,
 \quad p(x)=x s(x)-M(x),
\]
where all three are zero for $x<1$. Then $p$ is continuous, including
at integers, and $p(x)=x\eta(x)$ for $x\geq1$.
Let $v=\log2$ and define
\begin{equation}
 f(x)=x s(x)v+\sum_{x<n\leq2x}\mu(n)\log\frac{2x}{n}.
 \label{eq:ns68-forcing}
\end{equation}
In particular $f(x)=0$ for $x<1/2$.

\begin{proposition}[Exact first-interval response and renewal identity]
\label{prop:ns68-renewal}
For each integer $N\geq1$,
\begin{equation}
 f(N)=F_N(2N)=e_N(1/(2N)).
 \label{eq:ns68-canonical-sample}
\end{equation}
For every real $x>0$,
\begin{equation}
 f(x)=v p(x)+\int_x^{2x}\frac{M(u)}u\,du
     =p(2x)-(1-v)p(x)-\int_x^{2x}\frac{p(u)}u\,du.
 \label{eq:ns68-renewal-identity}
\end{equation}
The coefficients in the corresponding average have total mass
$(1-v)+v=1$, and both are nonnegative.
\end{proposition}
\begin{proof}
On $N<m<2N$, every proper divisor of $m$ is at most $N$, so
$q_N(m)=-\mu(m)$. The endpoint $m=2N$ has zero logarithmic weight.
Equation~\eqref{eq:ns67-divisor-response} therefore gives
\eqref{eq:ns68-canonical-sample}. It also holds for $N=1$ directly.
Partial summation gives
\[
 \sum_{x<n\leq2x}\mu(n)\log\frac{2x}{n}
       =\int_x^{2x}\frac{M(u)-M(x)}u\,du,
\]
including integer endpoints, since the coefficient vanishes at $2x$.
This proves the first identity. The continuous, piecewise differentiable
function $p$ satisfies $u p'(u)-p(u)=M(u)$ almost everywhere. Integrate
$p'(u)-p(u)/u$ from $x$ to $2x$ to get the second identity. These
identities also apply when the interval crosses $1$.
\end{proof}

\begin{proposition}[Weighted causal inversion, with no power gain]
\label{prop:ns68-weighted-inversion}
For $\theta>0$ define
\[
 q_\theta=(1-v)2^{-\theta}+\frac{1-2^{-\theta}}\theta<1.
\]
Write $P_\theta(t)=e^{-\theta t}p(e^t)$ and
$G_\theta(t)=e^{-\theta t}f(e^{t-v})$, both zero for $t<0$.
Then
\begin{equation}
 P_\theta(t)=(1-v)2^{-\theta}P_\theta(t-v)
       +\int_0^v e^{-\theta u}P_\theta(t-u)\,du+G_\theta(t).
 \label{eq:ns68-causal}
\end{equation}
For every $1\leq a\leq\infty$ and every $T>0$,
\begin{equation}
 (1-q_\theta)\|P_\theta\|_{L^a(0,T)}
 \leq\|G_\theta\|_{L^a(0,T)}
 \leq(1+q_\theta)\|P_\theta\|_{L^a(0,T)}.
 \label{eq:ns68-two-sided}
\end{equation}
These bounds also apply on $(0,\infty)$ whenever either side is
finite. In particular, a bound $|f(x)|\leq Cx^\theta$ on $x>0$
implies
\begin{equation}
 |p(x)|\leq\frac{2^{-\theta}C}{1-q_\theta}\,x^\theta.
 \label{eq:ns68-sup-transfer}
\end{equation}
No smaller power follows from this renewal identity for general forcing.
\end{proposition}
\begin{proof}
Substitute $x=e^{t-v}$ in~\eqref{eq:ns68-renewal-identity}. The
convolution kernel in~\eqref{eq:ns68-causal} is the positive measure
\[
 K_\theta=(1-v)2^{-\theta}\delta_v
              +e^{-\theta u}\mathbf1_{(0,v)}(u)\,du.
\]
Its mass is $q_\theta$. At $\theta=0$ the mass is one; for
$\theta>0$ all its mass is strictly damped at positive delays, so
$q_\theta<1$. Extend a restriction to $(0,T)$ by zero. Causality
makes its convolution agree with the original convolution on $(0,T)$.
Young's inequality and the triangle inequality give both sides of
\eqref{eq:ns68-two-sided}. The finite-interval estimate followed by
$T\to\infty$ justifies the assertion without assuming an unknown
infinite norm is finite. For the supremum bound,
$|G_\theta|\leq2^{-\theta}C$.

For an explicit scope check, take the nonarithmetic continuous function
$p_\theta(x)=(x^\theta-1)\mathbf1_{x\geq1}$ and define its forcing
by the renewal identity. For $x\geq1$ that forcing is exactly
\[
 f_\theta(x)=\left\{2^\theta-(1-v)
                  -\frac{2^\theta-1}{\theta}\right\}x^\theta
            =2^\theta(1-q_\theta)x^\theta.
\]
The coefficient is positive. Thus the inverse preserves this power,
which excludes a generic power improvement by positivity alone. This
is a test of the operator class, not a counterexample concerning the
actual M\"obius coefficients.
\end{proof}

\begin{proposition}[The single sample retains every nontrivial-zero obstruction]
\label{prop:ns68-scalar-criterion}
The following two statements are equivalent to RH:
\[
 \forall\varepsilon>0,\quad p(x)=O_\varepsilon(x^{1/2+\varepsilon}),
 \qquad
 \forall\varepsilon>0,\quad f(N)=O_\varepsilon(N^{1/2+\varepsilon}).
\]
The second assertion is only over integer $N$. No such bound is
proved here.
\end{proposition}
\begin{proof}
RH implies $M(x)=O_\delta(x^{1/2+\delta})$. The complete tail
formula for $\eta$ in NS-66 then bounds $p$ at the same power.
Conversely the Mellin identity
\begin{equation}
 \int_1^\infty p(x)x^{-s-1}\,dx
             =\frac1{(s-1)s\zeta(s)},\qquad \Re s>1,
 \label{eq:ns68-p-transform}
\end{equation}
continues analytically to $\Re s>1/2$ if all the stated power bounds
hold. A zero there would give a genuine pole, a contradiction.
The apparent singularity at $s=1$ is removable. The functional equation
then gives RH.

The renewal identity and~\eqref{eq:ns68-sup-transfer} transfer these
power bounds in both directions between $p$ and $f$ for real $x$.
For the discrete-to-continuous step, the exact divisor identity gives
$|s(x)|\leq1+1/x$ for $x\geq1$. Off the discrete breakpoints,
\[
 f'(x)=v s(x)+\frac{M(2x)-M(x)}x,
 \qquad |f'(x)|\leq2v+2<4.
\]
The sums defining $f$ match continuously at integers and half-integers:
a term entering at $2x=n$ has coefficient zero, and a term leaving
at $x=n$ is exactly replaced by the change in $xs(x)v$.
Therefore $f$ is globally $4$-Lipschitz on $[1,\infty)$.
Its integer values suffice for every displayed power bound. Compact
initial intervals change only the constant.
\end{proof}

For an independent analytic check, termwise integration on $\Re s>1$
gives the exact Mellin response
\begin{equation}
 \int_0^\infty f(x)x^{-s-1}\,dx
  =\frac{H(s)}{\zeta(s)},\qquad
 H(s)=\frac{v}{s-1}+\frac{2^s-1-sv}{s^2}.
 \label{eq:ns68-forcing-transform}
\end{equation}
The factor $H$ has no zero on $\Re s>0$ (and has a pole at $1$).
Indeed the unweighted probability kernel has Laplace transform
\[
 \widehat K(s)=(1-v)2^{-s}+\frac{1-2^{-s}}s,
 \qquad H(s)=\frac{2^s\{1-\widehat K(s)\}}{s(s-1)}.
\]
For $\sigma=\Re s>0$, positivity of that kernel gives
$|\widehat K(s)|\leq q_\sigma<1$. Hence the numerator never
vanishes in this half-plane. The apparent pole of $H/\zeta$ at
$s=1$ is removable, as in~\eqref{eq:ns68-p-transform}.
Thus no nontrivial-zeta-zero contribution is removed by this scalar
sampling kernel. This is the same weighted contraction in Mellin
coordinates, not an additional arithmetic estimate.

\begin{corollary}[Critical weighted total energy cannot be finite]
\label{cor:ns68-critical-energy}
Put $q=q_{1/2}$. For every $X\geq1$,
\begin{equation}
 (1-q)^2\int_1^X\eta(x)^2\,dx
 \leq\frac12\int_{1/2}^{X/2}\frac{f(x)^2}{x^2}\,dx
 \leq(1+q)^2\int_1^X\eta(x)^2\,dx.
 \label{eq:ns68-finite-energy-comparison}
\end{equation}
Both integrals diverge as $X\to\infty$, unconditionally. Finiteness
of this particular critical weighted total energy is therefore not a
possible sufficient estimate to pursue. This does not exclude a
logarithmic or subpolynomial growth bound for these integrals, and is
not a claim that the canonical norm $D_{N,1}$ diverges.
\end{corollary}
\begin{proof}
Use~\eqref{eq:ns68-two-sided} with $a=2$, $\theta=1/2$ and
$T=\log X$, and change variables. Here
$\|P_{1/2}\|^2=\int_1^X\eta(x)^2\,dx$ and
$\|G_{1/2}\|^2=\tfrac12\int_{1/2}^{X/2}f(x)^2x^{-2}\,dx$.

If $\int_1^\infty\eta(x)^2dx$ were finite, Cauchy--Schwarz would
make~\eqref{eq:ns68-p-transform} analytic on $\Re s>1/2$ and give,
for $s=1/2+\delta+i\gamma$,
\[
 \left|\int_1^\infty p(x)x^{-s-1}dx\right|
       \leq\frac{\|\eta\|_2}{\sqrt{2\delta}}.
\]
At any known critical-line zero $1/2+i\gamma$ its meromorphic
continuation has a pole of integer order at least one. Along the
horizontal approach that grows at least as $\delta^{-1}$, contradicting
the displayed bound. The positive lower factor in
\eqref{eq:ns68-finite-energy-comparison} transfers divergence to the
forcing energy. No assumption on simplicity or RH is needed.
\end{proof}

\paragraph{Finite verification.}
Exact Maxima checks retain integer and noninteger endpoints, the physical
canonical sample, the Mellin response and both energy primitives.
Arb replays at 256 and 384 bits verify the complete renewal identity
and finite energy comparison for
$N=16,64,512,4096,65536,1048576$, using every M\"obius coefficient
through $2N$. There is no omitted tail in these finite identities.
The critical kernel mass and supremum-transfer constant satisfy
\[
 .8027641<q_{1/2}<.8027642,\qquad
 3.58508<\frac{2^{-1/2}}{1-q_{1/2}}<3.58509.
\]
The archive matches $74$ cross-precision quantities. No growth bound
is inferred from these six values.

\paragraph{Stop condition.}
The averaging geometry yields an explicit inverse bound, not a new
arithmetic estimate. A square-root-plus-epsilon estimate for its signed
forcing is exactly strong enough for RH and is still absent. Bounding
all critical weighted energy by a finite constant is actually false.
The remaining permissible objectives are growth estimates with an
explicit rate; neither positivity of the renewal kernel nor the finite
checks supply them.
